\chapter{Coefficients binomiaux et formule du binôme}
\section*{Exercice 4.1 -- Pascal}
Depuis $(1,4,6,4,1)$, on obtient $(1,5,10,10,5,1)$. Sa somme vaut $32=2^5$. Les extrémités valent $1$ et chaque valeur intérieure somme deux voisines de la ligne précédente.

\section*{Exercice 4.2 -- Coefficient}
Le terme de degré deux est $\binom52(-2x)^2=40x^2$. Le signe est positif, car l'exposant de $-2$ est pair.

\section*{Exercice 4.3 -- Somme alternée}
Fixer un élément et l'ajouter ou le retirer à chaque partie échange taille paire et taille impaire. Cette opération est sa propre inverse ; elle donne une bijection entre les deux catégories, donc leurs nombres sont égaux.

\section*{Exercice 4.4 -- Élément distingué}
Compter les couples $(B,b)$ avec $B$ une partie et $b\in B$. Par taille de $B$, on obtient le membre de gauche. En choisissant d'abord $b$ parmi $n$ éléments puis les autres membres librement, on obtient $n2^{n-1}$.

\section*{Exercice 4.5 -- Somme de carrés}
Partager $2n$ objets en deux groupes de $n$. Pour choisir $n$ objets dont $k$ dans le premier groupe, on dispose de $\binom nk\binom n{n-k}=\binom nk^2$ choix. Les valeurs de $k$ partitionnent tous les choix.

\section*{Exercice 4.6 -- Binôme et matrices}
La distributivité donne $A^2+AB+BA+B^2$. La formule $A^2+2AB+B^2$ est valide si $AB=BA$. Sans commutativité, les deux termes croisés doivent rester séparés.
